%HOMEWORK template for M 325K
\documentclass{article}
\headheight=7pt         \topmargin=14pt
\textheight=574pt       \textwidth=445pt
\oddsidemargin=18pt     \evensidemargin=18pt

%Begin package import

\usepackage{amsmath, amssymb, amsthm}
\usepackage{mathtools}
\usepackage{pdfpages}
\usepackage{tikz-cd}

%End package import
%Begin custom commands

\newcommand{\C}{\mathbb{C}}
\newcommand{\R}{\mathbb{R}}
\newcommand{\Q}{\mathbb{Q}}
\newcommand{\Z}{\mathbb{Z}}
\newcommand{\N}{\mathbb{N}}
\newcommand{\F}{\mathbb{F}}
\newcommand{\abs}[1]{\left\lvert #1\right\rvert}
\newcommand{\RM}[1]{\mathrm{#1}}
\newcommand{\BF}[1]{\textbf{#1}}
\newcommand{\e}{\epsilon}
\newcommand{\ceil}[1]{\lceil{#1}\rceil}
\newcommand{\floor}[1]{\lfloor{#1}\rfloor}
\newcommand{\rarr}[0]{\rightarrow}
\newcommand{\IM}[1]{\text{Im}(#1)}
\newcommand{\Hom}[0]{\text{Hom}}
\newcommand{\Pow}[1]{\text{Pow}(#1)}
\newcommand{\angles}[1]{\langle #1 \rangle}
\newcommand{\ext}[2]{\text{Ext}_{#1}^{#2}}

%End custom commands
%Begin custom theorems

\newtheorem{innercustomgeneric}{\customgenericname}
\providecommand{\customgenericname}{}
\newcommand{\newcustomtheorem}[2]{%
\newenvironment{#1}[1]
{%
\renewcommand\customgenericname{#2}%
\renewcommand\theinnercustomgeneric{##1}%
\innercustomgeneric%
}
{\endinnercustomgeneric}
}

\newcustomtheorem{THRM}{Theorem}
\newcustomtheorem{LEM}{Lemma}
\newcustomtheorem{EX}{Exercise}
\newcustomtheorem{COR}{Corollary}
\newcustomtheorem{PROB}{Problem}
\newcustomtheorem{claim}{Claim}

\newenvironment{solution}
{\renewcommand\qedsymbol{$\blacksquare$}\begin{proof}[Solution]}
{\end{proof}}

\newenvironment{definition}
{\renewcommand\qedsymbol{$\blacksquare$}\begin{proof}[Definition]}
{\end{proof}}

%End custom theorems
\begin{document}

\title{Field Symmetries}
\author{Arham Lodha}%put your name here
\date{March 30, 2024}%enter the due date here
\maketitle

Take $L / K$ a finite extension. Goal is to bound $G := Gal(L / K)$.   

Consider another field $F$. Let $g: K \to F$ be a map of rings. 

Let $Ext_K^L(g) := \left\{ h: L \to F | h|_K = g \right\}$.  
\begin{PROB}{1}\label{Problem 1.}
    Show that if F = L, and g the original inclusion, then this Ext is just G. 
\end{PROB}
\begin{proof}
    $\forall h \in Ext_K^L(g)$, $h$ is a inclusion and thus a isomorphism of rings (more specifically automorphism). If $g$ is the original inclusion, it is a identity map. Thus $h|_K = id_K$. Thus by definition of Galois Group, $h \in G$. Thus $Ext_K^L(g) \subseteq G$ 
    $\forall h \in G$, $h|_K = id_K = g$, by definition of Galois Group. Thus $h \in Ext_K^L(g)$. Thus $G \subseteq Ext_K^L(g)$.
    Thus $Ext_K^L(g) = G$. 
\end{proof}

\bigskip

\begin{PROB}{2a}\label{Problem 2a.}
    Suppose $L = K[a]$, and $p(X) \in K[X]$ is the minimal monic poly. Construct a canonical bijection between Ext and the roots of $g(p)[X]$ in $F[X]$  (note g induces $L[X] \to F[X]$, a ring homo, just applying g to the coefficients and sending X to X).
\end{PROB}
\begin{proof}
    The canonical map between $\ext{K}{L}(g)$ and the roots of $g(p)(X)$ is the following $h \to h(a)$.

    Suppose $\forall h, k \in \ext{K}{L}(g)$, $h(a) = k(a) \implies \forall x + ya, h(x + ya) = g(x) + g(y)h(a) = g(x) + g(y)k(a) = k(x+ ya) \implies h = k$.
    
    $\forall r$ in the roots of $g(p)[X]$, there exists some function which sends $a \to r$. Thus the canonical map is surjective.
    
    Thus the map is bijection.  
\end{proof}

\begin{PROB}{2b}\label{Problem 2b.}
    Conclude $|Ext|$ is at most the degree of $L / K$, with equality iff $g(p)$ has all its roots in F and there are no double roots.
\end{PROB}
\begin{proof}
    Thus $|\ext{K}{L}|$ is at equal the number of unique roots of $g(p)$ in $F$ which is at most the number of unique roots of $p$ which is at most $\deg p = [L : K]$. Thus $|\ext{K}{L}| \leq [L : K]$.

    Suppose $|\ext{K}{L}| = [L : K]$. Then the number of unique roots of $g(p)$ in F is $[L : K]$. $\deg g(p) = \deg(p) = [L : K]$. Thus the number of unique roots of $g(p)$ in F is equal to the number of possible roots of $g(p)$ (or the degree). Thus every root is unique and $g(p)$ has no double root. Thus since the number of roots of $g(p)$ in $F$ is equal to the number of roots of $g(p)$, all roots are in $F$. Thus $g(p)$ has all roots in $F$ and no double roots. 
    
    Suppose $g(p)$ has all roots in $F$ and no double roots. The number of unique roots of $g(p)$ in $F$ is equal to $\deg g(p) = \deg p = [K : L]$. Since $|\ext{K}{L}|$ equals the number of number of unique roots of $g(p)$ in $F$, $|\ext{K}{L}| = [K : L]$.        
\end{proof}

\bigskip

\begin{PROB}{3}\label{Problem 3.}
    Suppose we have $K \subseteq L \subseteq M$, all fields.  Note there is a natural map $\alpha: \ext{K}{M}(g) \to \ext{K}{L}(g)$. Show that the fibre over q is exactly $\ext{L}{M}(q)$ . 
\end{PROB}
\begin{proof}
    $\alpha^*(q) := \left\{ h \in \ext{K}{M}(g) \mid \forall l \in L h(l) = q(l) \right\} = \left\{ h \in \ext{K}{M}(g) \mid h|_L = q \right\} = \ext{L}{M}(q)$, by the definition of ext.
\end{proof}

\bigskip

\begin{PROB}{4}\label{Problem 4.}
    Now show $\ext{K}{L}(g)$  has order at most degree L/K, with equality iff for every a in K, $g(p_a)$ has all its roots in F and these are all distinct, where $p_a$ means the minimal monic for a in K[X]. 
\end{PROB}
\begin{proof}
    Let $L_i := L_{i - 1}[\alpha]$ for some $\alpha \in L_i$. Let $L_0 = K$ and let there be some $N \in \N \ni \forall m \geq N, L_m = L$.    

    Base Case: $L_1 = K[\alpha]$. By problem 2b, $\ext{L_0}{L_1}(g) \leq \left[K[\alpha] : K\right]$ with the needed equality condition. 
    
    Inductive Hypothesis: Suppose $\ext{K}{L_{n - 1}}(g)$ satifies the theorem.
    
    Inductive Step: $L_{n} = L_{n - 1}[\alpha]$ for some $\alpha \in L_n$. 
    Take $r: \ext{K}{L_{n}}(g) \to \ext{K}{L_{n - 1}}(g)$. By the inductive hypothesis, $\ext{K}{L_{n - 1}}(g) \leq [L_{n - 1} : K]$. Furthermore from 3, we know $\forall q \in \ext{K}{L_{n - 1}}(g)$, $r^*(q) = \ext{L_{n - 1}}{L_n}(q)$   
    
    Let $p$ be the minimal polynomial of $\alpha$ over $L_{n - 1}$. 

    By 2b, $\ext{L_{n - 1}}{L_n}(q) \leq [L_{n - 1} : L_n]$ with equality when $q(p)$ has all its roots in $L_i$ with no double roots.

    Thus, $|\ext{K}{L_{n}}(g)| = \sum_{q \in \ext{K}{L_{n - 1}}(g)} |r^*(q)| \leq |\ext{K}{L_{n - 1}}(g)| |\ext{L_{n - 1}}{L_n}(q)| \leq [L_{n - 1} : K][L_{n} : L_{n - 1}] = [L_{n} : K]$. Thus $|\ext{K}{L_{n}}(g)| \leq [L_{n} : K]$
    
    
    

    Suppose, $\forall a \in L_n$, $p_a \in K[X]$, the minimum polynomial for $a$, $g(p_a)$ has all its roots in $F$ and they are all distinct. Since $K \subseteq L_{n - 1} \subseteq L_{n}$, then the minimal polynomial of $a$ in $L_{n - 1}$, $p'_a$, divides $p_a$. Since $g(p_a)$ has all its roots in $F$ and the roots are distinct, $q(p'_a)$ must also have all its roots and they must be distinct. Thus $\ext{L_{n - 1}}{L}(q) = [L_n : L_{n - 1}]$. Since $L_{n - 1} \subseteq L_n$, by induction hypothesis, $|\ext{K}{L_{n- 1}}| = [L_{n - 1} : K]$. Thus, $r$ is surjective since every element in $\ext{K}{L_{n- 1}}$ has the same fibre. Thus, since $\ext{K}{L}(g) = \sum_{q \in \ext{K}{L_{n - 1}}(g)} |r^*(q)| = \sum_{q \in \ext{K}{L_{n - 1}}(g)} [L_{n} : L_{n - 1}] = [L_{n - 1} : K] [L_n : L_{n - 1}] = [L : K]$
    
    
    Suppose, $\ext{K}{L_n}(g) = [L : K]$. $\forall a \in L_n$. $K \subseteq K[a] \subseteq L_n$. Thus, $\ext{K}{K[a]}(g) = [K[a] : K]$. Thus $g(p_a)$ factors completely over $F$ with distinct roots. 
    
    
    Thus the theorem is true for $\ext{K}{L}(g)$ by induction.   
\end{proof}

\bigskip

\begin{PROB}{5}\label{Problem 5.}
    Now go back to the original problem, with $G := Gal(L/K)$. Show $|G|$ is at most the degree with equality iff for every a in L, its monic min polynomial over K[X] has all its roots in L, and  these are all distinct. 
\end{PROB}
\begin{proof}
    Suppose $F = L$ and $g = id$. By 1, $G = \ext{K}{L}(id_K)$. By problem 4, $|G| = |\ext{K}{L}(id_K)| \leq [L : K]$ with equality iff for every a in L, its monic min polynomial over K[X] has all its roots in L, and  these are all distinct.   
\end{proof}

\begin{PROB}{6}\label{Problem 6.}
    Show in char 0 that any irreducible polynomial has distinct roots (in any field extension).
\end{PROB}
\begin{proof}
    Let $p$ be the given irreducible with a root $\alpha$. Assume the contrary, $p$ has a repeated root r. Thus $p'$, derivative of $p$ wrt x, also has a root at x, because the condition for a repeated root is $p(x) = 0$ and $p'(x) = 0$. Thus $\gcd(p, p') \neq f \in F$. However since $f$ is irreducible, its divisors are $f$ and units. Thus $\gcd(p, p') = p$. However in $char F \neq 0$, $f'$ is zero (if $f$) constant which contradicts our assumption of irreduciblility, or must have degree less than $f$. But then $f$ cannot divide $f'$. Thus a contradiction occurs. Thus $f$  has no repeated roots.     
\end{proof}


\end{document}